% \documentclass[a4paper,12pt]{article}{{{
\documentclass[a4paper,11pt]{article}
% -------------------------------------------------
% Set up the page margins.
% -------------------------------------------------
% \usepackage[text={6.5in,9.5in},centering]{geometry}
% \setlength{\topmargin}{-0.25in}
\usepackage[margin=1in]{geometry}

% -------------------------------------------------
% Load Packages here
% -------------------------------------------------
%\usepackage{hyperref}
\usepackage{graphicx,url,subfig}
\RequirePackage[OT1]{fontenc}
\RequirePackage{amsthm,amsmath,amssymb,amscd}
\RequirePackage[numbers]{natbib}
\RequirePackage[colorlinks,citecolor=blue,urlcolor=blue]{hyperref}
\usepackage{graphics}
\usepackage{enumerate}
\usepackage{datetime}
\usepackage{etex}
% \usepackage[notcite,notref,color]{showkeys}
\usepackage[normalem]{ulem} % either use this (simple) or
\usepackage{soul} % use this (many fancier options)
\makeatother
\numberwithin{equation}{section}
\allowdisplaybreaks
\usepackage{mathrsfs}


% -------------------------------------------------
% check references
% -------------------------------------------------
% \usepackage{refcheck}
% \usepackage[notref,notcite]{showkeys}
% \usepackage[notcite]{showkeys}
% \usepackage{showkeys}

% -------------------------------------------------
% Environments here
% -------------------------------------------------
\theoremstyle{plain}
\newtheorem{theorem}{Theorem}[section]
\newtheorem{corollary}[theorem]{Corollary}
\newtheorem{lemma}[theorem]{Lemma}
\newtheorem{proposition}[theorem]{Proposition}
\newtheorem{exercise}{Exercise}
\theoremstyle{definition}
\newtheorem{definition}[theorem]{Definition}
\newtheorem{hypothesis}[theorem]{Hypothesis}
\newtheorem{assumption}[theorem]{Assumption}

\newtheorem{remark}[theorem]{Remark}
\newtheorem{conjecture}[theorem]{Conjecture}
\newtheorem{example}[theorem]{Example}

% -------------------------------------------------
%  Define short hand symbols.
% -------------------------------------------------
\newcommand{\E}{\mathbb{E}}
\newcommand{\W}{\dot{W}}
\newcommand{\B}{\textbf{B}}
\newcommand{\D}{\mathbb{D}}
\newcommand{\ud}{\ensuremath{\mathrm{d} }}
\newcommand{\Ceil}[1]{\left\lceil #1 \right\rceil}
\newcommand{\Floor}[1]{\left\lfloor #1 \right\rfloor}
\newcommand{\sgn}{\text{sgn}}
\newcommand{\Lad}{\text{L}_{\text{ad}}^2}
\newcommand{\SI}[1]{\mathcal{I}\left[#1 \right]}
\newcommand{\SIB}[2]{\mathcal{I}_{#2}\left[#1 \right]}
\newcommand{\Indt}[1]{1_{\left\{#1 \right\} }}
\newcommand{\LadInPrd}[1]{\left\langle #1 \right\rangle_{\text{L}_\text{ad}^2}}
\newcommand{\LadNorm}[1]{\left|\left|  #1 \right|\right|_{\text{L}_\text{ad}^2}}
\newcommand{\Norm}[1]{\left\|  #1   \right\|}
\newcommand{\Ito}{It\^{o} }
\newcommand{\Itos}{It\^{o}'s }
\newcommand{\spt}[1]{\text{supp}\left(#1\right)}
\newcommand{\InPrd}[1]{\left\langle #1 \right\rangle}
\newcommand{\mr}{\textbf{r}}
\newcommand{\Ei}{\text{Ei}}
\newcommand{\arctanh}{\operatorname{arctanh}}
\newcommand{\ind}[1]{\mathbb{I}_{\left\{ {#1} \right\} }}

\DeclareMathOperator{\esssup}{\ensuremath{ess\,sup}}

\newcommand{\steps}[1]{\vskip 0.3cm \textbf{#1}}
\newcommand{\calB}{\mathcal{B}}
\newcommand{\calD}{\mathcal{D}}
\newcommand{\calE}{\mathcal{E}}
\newcommand{\calF}{\mathcal{F}}
\newcommand{\calG}{\mathcal{G}}
\newcommand{\calK}{\mathcal{K}}
\newcommand{\calH}{\mathcal{H}}
\newcommand{\calI}{\mathcal{I}}
\newcommand{\calL}{\mathcal{L}}
\newcommand{\calM}{\mathcal{M}}
\newcommand{\calN}{\mathcal{N}}
\newcommand{\calO}{\mathcal{O}}
\newcommand{\calT}{\mathcal{T}}
\newcommand{\calP}{\mathcal{P}}
\newcommand{\calR}{\mathcal{R}}
\newcommand{\calS}{\mathcal{S}}
\newcommand{\bbC}{\mathbb{C}}
\newcommand{\bbN}{\mathbb{N}}
\newcommand{\bbP}{\mathbb{P}}
\newcommand{\bbZ}{\mathbb{Z}}
\newcommand{\myVec}[1]{\overrightarrow{#1}}
\newcommand{\sincos}{\begin{array}{c} \cos \\ \sin \end{array}\!\!}
\newcommand{\CvBc}[1]{\left\{\:#1\:\right\}}
\newcommand*{\one}{ {{\rm 1\mkern-1.5mu}\!{\rm I} }}
\newcommand{\RR}{\mathbb{R}}
\newcommand{\R}{\mathbb{R}}
\newcommand{\myEnd}{\hfill$\square$}
\newcommand{\ds}{\displaystyle}
\newcommand{\Shi}{\text{Shi}}
\newcommand{\Chi}{\text{Chi}}
\newcommand{\Daw}{\ensuremath{\mathrm{daw} }}
\newcommand{\Erf}{\ensuremath{\mathrm{erf} }}
\newcommand{\Erfi}{\ensuremath{\mathrm{erfi} }}
\newcommand{\Erfc}{\ensuremath{\mathrm{erfc} }}
\newcommand{\He}{\ensuremath{\mathrm{He} }}

\def\crtL{\theta}

\newcommand{\conRef}[1]{(#1)}

\DeclareMathOperator{\Lip}{\mathit{L}}
\DeclareMathOperator{\LIP}{Lip}
\DeclareMathOperator{\lip}{\mathit{l}}

\DeclareMathOperator{\Vip}{\overline{\varsigma}}
\DeclareMathOperator{\vip}{\underline{\varsigma}}
\DeclareMathOperator{\vv}{\varsigma}

% % Set up mdframe for questions.
\newcounter{NQ}
\newcounter{LQ}
\newcounter{ZQ}
\usepackage{xcolor}
\usepackage[framemethod=tikz]{mdframed}
\mdfsetup{%
	middlelinecolor=lightgray,
	middlelinewidth=2pt,
	backgroundcolor=red!20,
	roundcorner=10pt}
\newenvironment{NickQuestion}[1][] %
{\refstepcounter{NQ}
	\begin{mdframed}[backgroundcolor=lightgray]\textbf{Nick's Question \theNQ \: #1~~}~~~\noindent}%
	{\end{mdframed}}
\newenvironment{LeQuestion}[1][] %
{\refstepcounter{LQ}
	\begin{mdframed}[backgroundcolor=red!10]\textbf{Le's Question \theLQ \: #1~~}~~~\noindent}%
	{\end{mdframed}}
\newenvironment{ZhijianQuestion}[1][] %
{\refstepcounter{ZQ}
	\begin{mdframed}[backgroundcolor=blue!10]\textbf{Zhijian's Question \theZQ \: #1~~}~~~\noindent}%
	{\end{mdframed}}
\numberwithin{NQ}{section}
\numberwithin{LQ}{section}
\numberwithin{ZQ}{section}

\usepackage[notref,notcite]{showkeys}

\title{Nick Eisenberg's Meeting Notes}
\author{Le Chen, Nick Eisenberg, Zhijian Wu}
\date{October 2019}

\usepackage{natbib}
\usepackage{graphicx}
% }}}

\begin{document}

\begin{NickQuestion}[Inequality After Equation (B.11)]
	On the first line of the proof of Theorem 1.3 of \cite{BC2017} references the bounds from Lemma B.3. My question is about the first inequality right after equation (B.11). I will underline the parts of the inequality that this question pertains to. Here is the inequality:
		\begin{align*}
			| \mathcal{F} & (g_{ \textbf{t}, t+h, x }^{(n)}  - g_{ \textbf{t}, t, x }^{(n)} ) (\xi_1 , \cdots , \xi_n)|^2 
		\\& \le \lambda^{2n} \underbrace{ J_+^2(t,x) }_{(1)} \prod_{k=1}^{n-1} \exp\left( - \frac{ t_{\rho(k+1)}  - t_{\rho(k)} }{ t_{\rho(k)}t_{\rho(k+1)} } \left| \sum_{j=1}^k t_{\rho(j)}\xi_{\rho(j)} \right|^2  \right)
		\\&  \left| \exp\left( -\frac{1}{2} \frac{ t+h  - t_{\rho(n)} }{ t_{\rho(n)} (t+h) } \left| \sum_{j=1}^n t_j\xi_j \right|^2  \right)  -  \exp\left( -\frac{1}{2} \frac{ t - t_{\rho(n)} }{ t_{\rho(n)} t} \left| \sum_{j=1}^n t_j\xi_j \right|^2  \right)   \right|^2
		\end{align*}
	and it says that this follows from equation (3.5) which says that 
		\begin{align}
			 \mathcal{F} & g_{ \textbf{t}, t, x }^{(n)} (\xi_1 , \cdots , \xi_n) = \lambda^n \prod_{k=1}^n \exp\left( -\frac{1}{2}\frac{ t_{\rho(k+1)}  - t_{\rho(k)} }{ t_{\rho(k)}t_{\rho(k+1)} } \left| \sum_{j=1}^k t_{\rho(j)}\xi_{\rho(j)} \right|^2 \right)
			 \\& \times \exp \left( -\frac{i}{t} \left( \sum_{j=1}^n t_j \xi_j \right) \cdot x \right) \int_{\R^d} \exp \left( -i \left[ \sum_{j=1}^n \left( 1- \frac{t_j}{t} \right) \xi_j  \right] \cdot x_0 \right) \underbrace{ G(t,x-x_0) }_{(2)} u_0(\ud x_0).
		\end{align}
		
I do not understand why we can factor out $J_+(t,x)$ which is in the (1) underbrace above. If we calculate $	| \mathcal{F} (g_{ \textbf{t}, t+h, x }^{(n)}  - g_{ \textbf{t}, t, x }^{(n)} ) (\xi_1 , \cdots , \xi_n)|^2 $ by using equation (3.5), then because of the (2) underbrace above we would have that 
	\begin{align*}
			| \mathcal{F} & (g_{ \textbf{t}, t+h, x }^{(n)}  - g_{ \textbf{t}, t, x }^{(n)} ) (\xi_1 , \cdots , \xi_n)|^2 = \Bigg| \cdots \int_{R^d} \exp \left( -i \left[ \sum_{j=1}^n \left( 1- \frac{t_j}{t+h} \right) \xi_j  \right] \cdot x_0 \right)G(t+h,x-x_0)u_0(\ud x_0) 
			\\& - \cdots \int_{R^d} \exp \left( -i \left[ \sum_{j=1}^n \left( 1- \frac{t_j}{t} \right) \xi_j  \right] \cdot x_0 \right)G(t,x-x_0)u_0(\ud x_0) \Bigg|^2
	\end{align*}
and so I do not see why we can factor this out as $J_+(t,x)$ because the integrand of the first term has $(t+h)$ while the second only has $(t)$. 

\noindent\textbf{Relation to Raluca's paper \cite{BS2017}:} The corresponding result is found in Lemma 7.2 \cite{BS2017}. The proof is exactly the same up until Raluca's equation (80). However in equation (81), you can see that Raluca took a different approach to bound $\psi_{t,h,x}^{(n)}(\textbf{t,t})$ and in return the $A_n(t,h,x)$ function.
\end{NickQuestion}
	

\addcontentsline{toc}{section}{Bibliography}
\begin{thebibliography}{999}
	\bibitem{BC2017}
	Balan, R. and Le, C. (2017)
	\newblock Parabolic Anderson Model with Space-Time Homogeneous Gaussian Noise and Rough Initial Condition
	\newblock {\em  J Theor. Probab.} 
	
	\bibitem{BS2017}
	Balan, R. and Song, J. (2017)
	\newblock Hyperbolic Anderson Model withspace-time homogeneous Gaussian noise
	\newblock {\em arXiv:1602.07004v2} 
\end{thebibliography}

\end{document}